\chapter{Legal Documentation}\label{ch:fm:legal-documentation}

When Lehman Brothers filed for bankruptcy in September 2008, Metavante Corporation had an interest rate swap on \$600 million of notional with
one of its subsidiaries, and the swap had moved in the subsidiary's favour. Metavante neither terminated nor paid: it relied on a clause of the
standard master agreement that makes each payment conditional on no \omterm{def:fm:legal-documentation:eod}{event of default} having occurred with respect to the other party. A year
later, in September 2009, the bankruptcy court in New York ruled from the bench that ``riding the market for the period of one year, while taking
no action whatsoever'' was contrary to the spirit of the Bankruptcy Code, that Metavante had waived its right to terminate by not exercising it
promptly, and that it had to pay. In April 2012 the Court of Appeal in London read the same clause in the Lehman European administration and held
that, under English law, a non-defaulting party's payment obligation stays suspended for as long as the default continues. A few lines of a
standard document decided who owed what to whom, and two courts read them differently.

\section{The master agreement and why it exists}

A firm trading over the counter signs one master agreement with each counterparty, and every trade between them becomes a confirmation under it.
Book 2 introduced the ISDA master agreement and the repo market's GMRA; Book 6 the close-out netting they make possible. The design point is that
all the trades under one master form a single agreement, so that on a default they are terminated together and netted to one amount, rather
than cherry-picked by an administrator who performs the profitable ones and disclaims the rest.

The agreement's printed form is standard; the negotiation is in its schedule, which elects and amends the printed terms, and in the collateral
annex. What a trading firm's legal and treasury staff negotiate is a short list: which events let each side terminate, how much collateral moves
and when, what collateral is eligible, and what happens on the first day something goes wrong (\cref{fig:fm:legal-documentation:stack}).

\begin{omfigure}
\begin{tikzpicture}[font=\footnotesize,>=stealth,box/.style={draw,rounded corners,align=center,minimum height=0.7cm,text width=3.3cm}]
\node[box,fill=omDef!15] (m) at (0,2.4) {master agreement\\{\scriptsize printed form}};
\node[box,fill=omDef!25] (s) at (0,1.2) {schedule\\{\scriptsize elections, triggers}};
\node[box,fill=omProp!15] (c) at (-4.2,0) {confirmations\\{\scriptsize one per trade}};
\node[box,fill=omThm!15] (a) at (4.2,0) {collateral annex\\{\scriptsize threshold, MTA, eligible assets}};
\node[box,fill=gray!15,text width=2.8cm] (pb) at (-4.9,2.4) {prime-brokerage agreement};
\node[box,fill=gray!15,text width=2.8cm] (cl) at (4.9,2.4) {clearing and give-up agreements};
\draw[->] (m) -- (s); \draw[->] (s) -- (c); \draw[->] (s) -- (a);
\draw[dashed] (pb) -- (m.west); \draw[dashed] (cl) -- (m);
\node[font=\scriptsize,align=center] at (0,-0.4) {one agreement:\\terminated and netted together};
\end{tikzpicture}

\omcaption{The documents of an over-the-counter relationship: the master agreement and its schedule, the confirmations that become part of it, and the
collateral annex; the prime-brokerage, clearing and give-up agreements sit beside it and may be linked by cross-default. Schematic.}\label{fig:fm:legal-documentation:stack}
\end{omfigure}

\section{Events of default, termination events and close-out}

\begin{definition}[Event of default, termination event]\label{def:fm:legal-documentation:eod}
An \emph{event of default}\index{event of default} is an event attributable to one party, such as failure to pay or deliver, breach, insolvency or
default under other agreements, that allows the other party to terminate all transactions under the master agreement. A \emph{termination
event}\index{termination event} is an event that is no party's fault, such as illegality or a change of tax law, and that allows termination of the
affected transactions.
\end{definition}

\begin{definition}[Additional termination event, NAV trigger, cross-default clause]\label{def:fm:legal-documentation:ate}
An \emph{additional termination event}\index{additional termination event} is a \omterm{def:fm:legal-documentation:eod}{termination event} the parties add in the schedule. A \emph{NAV
trigger}\index{NAV trigger} is an additional \omterm{def:fm:legal-documentation:eod}{termination event} against a fund: a fall in its net asset value by more than agreed percentages over
one, three or twelve months, or below a floor. A \emph{cross-default clause}\index{cross-default clause} makes a party's default under other
agreements for borrowed money, above a threshold amount, an \omterm{def:fm:legal-documentation:eod}{event of default} under this one.
\end{definition}

The master agreement's section 2(a)(iii), as the Court of Appeal quoted it, makes each payment obligation subject to ``the condition precedent that
no Event of Default or Potential Event of Default with respect to the other party has occurred and is continuing''. The clause lets the
non-defaulting party stop paying while it decides whether to terminate. How long it may wait is where the two courts of the hook diverged.

\begin{definition}[Flawed-asset clause]\label{def:fm:legal-documentation:flawed}
A \emph{flawed-asset clause}\index{flawed-asset clause} makes a party's right to receive a payment conditional on an event that its own default
removes, so that the asset it holds (the right to be paid) is flawed from the start rather than taken away on insolvency; section 2(a)(iii) is the
standard example.
\end{definition}

The English court held that the underlying debt is undisturbed and only the payment obligation is suspended, that the suspension has no end date and
continues until the default is cured, and that the clause does not offend the principle against depriving an insolvent estate of its assets. The US
bankruptcy court, reading the Bankruptcy Code's safe harbours as protecting the right to terminate and net rather than the right to wait, held that
waiting a year had waived the right. A firm facing a defaulting counterparty in the money to the defaulter therefore has a decision to make, promptly,
whose answer depends on the governing law and the forum.

\begin{dated}{2026-09}{Stays on default rights and cleared-swap collateral}\label{dat:fm:legal-documentation:stays}
\textbf{US stay rules.} Under 12 CFR 252.83 and 252.84, the qualified financial contracts of covered US banking entities must provide that the
transfer and stay provisions of the US special resolution regimes are effective as if the contracts were governed by US law, and may not permit the
exercise of default rights related to an affiliate of the direct party entering insolvency proceedings. Counterparties of such banks accept these
terms in their agreements or through industry protocols. \textbf{Cleared swaps.} Under 17 CFR 22.2, a futures commission merchant must treat its
customers' cleared swaps and their collateral as belonging to the customers and segregate the collateral.
\end{dated}

\section{Collateral annexes}

The collateral annex turns the master agreement's credit exposure into daily transfers. Book 6 defined its terms: the threshold below which no
collateral is due, the minimum transfer amount below which no transfer is made, and the independent amount posted regardless of exposure. The
annex adds a rounding convention, a list of eligible collateral and the haircut on each.

\begin{example}[A call under an annex]\label{ex:fm:legal-documentation:call}
An annex has a threshold of \$2 million, a minimum transfer amount of \$0.5 million and a rounding of \$0.1 million; the secured party holds \$5
million of cash. At an exposure of \$10.37 million the credit support amount is $10.37-2=8.37$ and the call is $8.37-5=3.37$, rounded up to \$3.4
million. At \$6.8 million the difference, a return of \$0.2 million, is below the minimum transfer amount: nothing moves. Under dealer A's annex of
the tutorial, with an independent amount of \$5 million and a 2\% haircut on bonds, an exposure of \$4.2 million against \$3 million of cash and \$5
million of bonds calls $4.2+5-3-4.9=\$1.3$ million.
\end{example}

\omcode{code/firm/docterms/firm_docterms.py}{25}{48}{A collateral annex as data, and its daily call: threshold, independent amount, haircut value of
collateral held, minimum transfer amount and rounding.}\label{lst:fm:legal-documentation:annex}

\section{Prime-brokerage and clearing agreements}

\begin{definition}[Prime-brokerage agreement]\label{def:fm:legal-documentation:pba}
A \emph{prime-brokerage agreement}\index{prime-brokerage agreement} sets the terms on which a prime broker holds, finances and lends securities to a
client: margin (chapter 14's \omterm{def:fm:treasury-and-funding:house}{house margin} and any lock-up), the broker's rights of rehypothecation (Book 1) over the client's assets, its rights to
close out positions and set off amounts on default, and the client's events of default and \omterm{def:fm:legal-documentation:eod}{termination events}, including \omterm{def:fm:legal-documentation:ate}{NAV triggers}.
\end{definition}

\omterm{def:fm:legal-documentation:pba}{Prime-brokerage agreements} are where \omterm{def:fm:legal-documentation:ate}{NAV triggers} matter most, because they sit beside the broker's discretion over margin: a fund in a drawdown
may trip a trigger at the moment its brokers raise margin, and a trip at one counterparty can become a default at others through cross-default.
Clearing agreements with a futures commission merchant (Book 1) and give-up agreements between executing and clearing brokers complete the set;
their terms are shorter because the clearing house's rulebook carries most of the weight.

\section{Tutorial: the trigger}\label{tut:fm:legal-documentation:tut}

\begin{tutorial}
\textbf{Goal.} Load a fund's agreement terms, run a year of NAV and exposures, compute the daily collateral calls, find the first day each \omterm{def:fm:legal-documentation:ate}{NAV
trigger} trips, and close out the first counterparty to terminate. \textbf{End state:} the NAV chart with the trips
(\cref{fig:fm:legal-documentation:nav}) and the close-out table (\cref{tab:fm:legal-documentation:closeout}).
\begin{tutsteps}
\item \textbf{The fund.} NAV starts at \$400 million; a quiet half-year, a three-month drawdown, a partial recovery
(\texttt{fm\_docs.nav\_path}, synthetic).
\item \textbf{The terms.} Two prime brokers and two dealers have \omterm{def:fm:legal-documentation:ate}{NAV triggers} over 21, 63 and 252 business days, one with a floor of \$280
million; three dealers have collateral annexes (\texttt{fm\_docs.TRIGGERS}, \texttt{ANNEXES}; illustrative terms).
\item \textbf{The monitor.} \texttt{firm.docterms.first\_trip} finds each trigger's first trip; \texttt{fm\_docs.calls} runs the annexes daily.
\item \textbf{The close-out.} The first dealer to trip terminates: \texttt{firm.docterms.close\_out} nets its swaps, the collateral it holds and a
repo, with and without set-off.
\end{tutsteps}
\end{tutorial}

\begin{omfigure}
\begin{tikzpicture}
\begin{axis}[width=10.5cm,height=5.2cm,xlabel={\small business day},ylabel={\small NAV (\$ million)},xmin=0,xmax=251,ymin=220,ymax=500,
  tick label style={font=\footnotesize},grid=major,grid style={gray!20},table/col sep=comma]
\addplot[omDef,thick] table[x=day,y=nav]{figdata/desk/18-legal-documentation/nav.csv};
\draw[gray,dashed] (axis cs:0,280) -- (axis cs:251,280) node[pos=0.35,below,font=\scriptsize] {dealer A floor};
\draw[omProp,dotted,thick] (axis cs:137,220) -- (axis cs:137,470);
\draw[omProp,dotted,thick] (axis cs:150,220) -- (axis cs:150,470);
\draw[omProp,dotted,thick] (axis cs:177,220) -- (axis cs:177,470);
\draw[omProp,dotted,thick] (axis cs:180,220) -- (axis cs:180,470);
\node[font=\scriptsize,anchor=south east] at (axis cs:137,472) {B};
\node[font=\scriptsize,anchor=south west] at (axis cs:150,472) {PB1};
\node[font=\scriptsize,anchor=south east] at (axis cs:177,438) {PB2};
\node[font=\scriptsize,anchor=south west] at (axis cs:180,472) {A};
\end{axis}
\end{tikzpicture}

\omcaption{The fund's NAV over the year and the first trip of each counterparty's \omterm{def:fm:legal-documentation:ate}{NAV trigger}: dealer B (a 12\% fall in 21 days) on day 137,
prime broker 1 (15\% in 21 days) on day 150, prime broker 2 (30\% in 63 days) on day 177 and dealer A (the \$280 million floor) on day 180 (labels B, PB1, PB2, A). Data:
\texttt{fm\_docs.nav\_path}, \texttt{fm\_docs.trips}.}\label{fig:fm:legal-documentation:nav}
\end{omfigure}

The fund peaks at \$447.7 million on day 45 and falls to \$247.7 million on day 204, 44.7\% below the peak. The triggers trip in the order of their
tightness, not of the fund's troubles: dealer B's 12\% in a month trips on day 137 when the fund is still at \$369.3 million, 12.4\% below its level
a month earlier; the prime brokers follow within six weeks; dealer A's floor trips last, on day 180 at \$275.2 million. Four counterparties can
terminate within 43 business days, which is the risk the treasurer of chapter 14 must plan for: the terms of the documents decide when the funding
stops.

\begin{omfigure}
\begin{tikzpicture}
\begin{axis}[width=10.5cm,height=5cm,xlabel={\small business day},ylabel={\small \$ million},xmin=0,xmax=251,ymin=-3,ymax=20,
  tick label style={font=\footnotesize},grid=major,grid style={gray!20},table/col sep=comma,
  legend cell align=left,legend style={font=\scriptsize,at={(0.5,-0.3)},anchor=north,/tikz/every even column/.append style={column sep=8pt}},legend columns=2]
\addplot[omDef,thick,const plot] table[x=day,y=held_A]{figdata/desk/18-legal-documentation/collateral.csv};
\addplot[omDef!50,thin] table[x=day,y=exp_A]{figdata/desk/18-legal-documentation/collateral.csv};
\addplot[omThm,thick,const plot] table[x=day,y=held_B]{figdata/desk/18-legal-documentation/collateral.csv};
\addplot[omThm!50,thin] table[x=day,y=exp_B]{figdata/desk/18-legal-documentation/collateral.csv};
\legend{held by dealer A,exposure of dealer A,held by dealer B,exposure of dealer B}
\end{axis}
\end{tikzpicture}

\omcaption{Collateral the fund has posted to two dealers against their exposure to it. Dealer A's annex adds a \$5 million independent amount, so it
holds more than its exposure throughout; dealer B's \$2 million threshold means it holds less. Dealer C's \$10 million threshold is never reached.
Data: \texttt{fm\_docs.calls}.}\label{fig:fm:legal-documentation:collateral}
\end{omfigure}

Over the year the annexes move collateral on 87 days for dealer A (44 calls, 43 returns) and 39 for dealer B, and never for dealer C, whose \$10
million threshold is above its largest exposure. Dealer B terminates on day 137.

\begin{table}[ht]
\centering\small
\begin{tabular}{lr}
\toprule
dealer B, day 137 & \$ million \\
\midrule
swaps: exposure at mid & 2.84 \\
swaps: replacement cost added by the dealer & 3.00 \\
less collateral the dealer holds & $-$1.20 \\
swaps balance, owed by the fund & 4.64 \\
repo: excess collateral the dealer holds, owed to the fund & $-$6.00 \\
\midrule
net with set-off (owed to the fund) & $-$1.36 \\
without set-off & fund pays 4.64; claims 6.00 \\
\bottomrule
\end{tabular}
\caption{Close-out of dealer B's agreements with the fund on its trigger date. Data: \texttt{fm\_docs.close\_out\_first}.}\label{tab:fm:legal-documentation:closeout}
\end{table}

With a set-off clause across the dealer's agreements, the fund's debt on the swaps and the dealer's debt on the repo net to \$1.36 million owed to
the fund. Without one, the fund must pay \$4.64 million under the swaps and pursue \$6 million under the repo separately, which is the same money
only if the dealer pays; the exposure to the dealer's credit is the gross \$6 million, not the net.

\section{Exchange memberships and give-ups}

Exchange membership (Book 1) is itself a contract: the rulebook binds the member, its traders and its systems, and the clearing agreement with a
clearing member guarantees the member's trades to the clearing house. A firm that executes through several brokers and clears through one signs
give-up agreements (Book 2's give-up lines) that let the executing brokers pass its trades to its clearer, within limits the clearer sets. The
documents are shorter than a master agreement and change less, but the limits in them are operational: a give-up line that is full stops trading.

\begin{method}[Negotiating a fund's documents]\label{met:fm:legal-documentation:negotiate}
\begin{enumerate}
\item List every agreement and, for each, its events of default, \omterm{def:fm:legal-documentation:eod}{termination events} and triggers, with their thresholds and cure periods.
\item Simulate the fund's NAV in a bad year and find the order in which triggers trip; widen the tightest.
\item Align thresholds, minimum transfer amounts and eligible collateral with treasury's forecast (chapter 14).
\item Negotiate set-off across agreements with the same counterparty group, and cross-default thresholds high enough not to cascade.
\item Keep the terms as data, and monitor them daily.
\end{enumerate}
\end{method}

\section{Build: agreement terms as data}\label{bld:fm:legal-documentation:docterms}

\begin{build}
\bfield{Purpose} No honest build: the subject is legal documentation. The analytical tool keeps agreement terms as data and computes what they
imply day by day.
\bfield{Interface} \texttt{firm.docterms}: \texttt{Annex}, \texttt{held\_value}, \texttt{call}; \texttt{Trigger}, \texttt{first\_trip};
\texttt{close\_out}.
\bfield{Rules} Ineligible collateral counts for nothing; no transfer below the minimum transfer amount; calls round up and returns round down;
triggers compare each day's NAV with the NAV a window earlier and with the floor.
\bfield{Acceptance tests} \texttt{code/firm/docterms/tests/}: annex calls on hand numbers; trigger trips by fall and by floor; close-out with and
without set-off.
\bfield{Stretch} Cure periods and notice periods; triggers on the NAV at month ends; a cross-default cascade across the fund's counterparties.
\end{build}

\begin{omsources}
\item Lomas v JFB Firth Rixson Inc [2012] EWCA Civ 419.
\item Davis Polk, ``Lehman Bankruptcy Court Rules \dots'', client memorandum, 29 September 2009, on In re Lehman Brothers Holdings Inc.
(Metavante).
\item 12 CFR 252.83 and 252.84; 17 CFR 22.2.
\end{omsources}

\section{Exercises}

\begin{exercise}[$\star$]\label{exo:fm:legal-documentation:1}
Under the annex of \cref{ex:fm:legal-documentation:call}, what moves at exposures of \$2.3 million and \$10.37 million with no collateral held?
\end{exercise}

\begin{exercise}[$\star$]\label{exo:fm:legal-documentation:2}
What is the difference between an \omterm{def:fm:legal-documentation:eod}{event of default} and a \omterm{def:fm:legal-documentation:eod}{termination event}? Which is a \omterm{def:fm:legal-documentation:ate}{NAV trigger}?
\end{exercise}

\begin{exercise}[$\star$]\label{exo:fm:legal-documentation:3}
On which days do the tutorial's four triggers trip, and which trips first?
\end{exercise}

\begin{exercise}[$\star\star$]\label{exo:fm:legal-documentation:4}
Why did Metavante neither pay nor terminate, and what did the New York court decide?
\end{exercise}

\begin{exercise}[$\star\star$]\label{exo:fm:legal-documentation:5}
What did the English Court of Appeal decide about section 2(a)(iii), and how does it differ from the New York ruling?
\end{exercise}

\begin{exercise}[$\star\star$]\label{exo:fm:legal-documentation:6}
Why does dealer A hold more collateral than its exposure, and dealer C none?
\end{exercise}

\begin{exercise}[$\star\star\star$]\label{exo:fm:legal-documentation:7}
\emph{Coding.} Tighten prime broker 2's 63-day trigger from 30\% to 20\% and rerun \texttt{fm\_docs.trips}. Which trips first now?
\end{exercise}

\begin{exercise}[$\star\star\star$]\label{exo:fm:legal-documentation:8}
\emph{Find the flaw.} ``Our \omterm{def:fm:legal-documentation:ate}{NAV triggers} are generous: none trips unless we lose 30\%.''
\end{exercise}

\section{Problem: The Trigger}

\begin{problem}[{Weekend problem --- the trigger}]\label{pb:fm:legal-documentation:1}
A fund's general counsel must review its agreements with two prime brokers and three dealers before a volatile year.

\medskip\noindent\textbf{Part I --- The documents.}
\begin{enumerate}
\item Why is a master agreement a single agreement, and what does that give a non-defaulting party?
\item Define an \omterm{def:fm:legal-documentation:eod}{event of default}, a \omterm{def:fm:legal-documentation:eod}{termination event}, an \omterm{def:fm:legal-documentation:ate}{additional termination event} and a \omterm{def:fm:legal-documentation:ate}{NAV trigger}.
\item Define a \omterm{def:fm:legal-documentation:ate}{cross-default clause} and a \omterm{def:fm:legal-documentation:flawed}{flawed-asset clause}.
\item What does section 2(a)(iii) say, as the Court of Appeal quoted it?
\end{enumerate}

\medskip\noindent\textbf{Part II --- The cases.}
\begin{enumerate}[resume]
\item Summarise Metavante and its outcome.
\item Summarise Lomas and its holding on the suspension.
\item What should a non-defaulting party decide, and how fast, after a counterparty's default?
\item What do the US stay rules change for counterparties of large US banks?
\end{enumerate}

\medskip\noindent\textbf{Part III --- The year.}
\begin{enumerate}[resume]
\item Give the NAV's peak, trough and drawdown.
\item Give each trigger's first trip and its reason.
\item Give the collateral moves under each annex, and explain dealer C.
\item Compute the call of \cref{ex:fm:legal-documentation:call}.
\item Give dealer B's close-out with and without set-off.
\end{enumerate}

\medskip\noindent\textbf{Part IV --- The review.}
\begin{enumerate}[resume]
\item Define a \omterm{def:fm:legal-documentation:pba}{prime-brokerage agreement} and name the terms that interact with chapter 14's margin.
\item Which triggers would you renegotiate first, and why?
\item What does cross-default do to the order of trips?
\item Why does set-off matter more when the counterparty is weak?
\item What do give-up agreements add to the picture?
\item State the \emph{named result}: the first day a declining NAV trips an \omterm{def:fm:legal-documentation:ate}{additional termination event} under each counterparty's terms, and the
close-out amount after set-off.
\item In two sentences, write the review's recommendation.
\end{enumerate}
\end{problem}

\section{Interview questions}

\begin{interviewq}[$\star$ \iqroles{trader}]\label{iq:fm:legal-documentation:1}
What is a minimum transfer amount, and why does an annex have one?
\end{interviewq}

\begin{interviewq}[$\star$ \iqroles{risk}]\label{iq:fm:legal-documentation:2}
What is a \omterm{def:fm:legal-documentation:ate}{NAV trigger}, and why do counterparties ask funds for one?
\end{interviewq}

\begin{interviewq}[$\star\star$ \iqroles{developer}]\label{iq:fm:legal-documentation:3}
Implement a collateral call with a threshold, an independent amount, haircuts, a minimum transfer amount and rounding. What are the edge cases?
\end{interviewq}

\begin{interviewq}[$\star\star$ \iqroles{risk}]\label{iq:fm:legal-documentation:4}
Your counterparty has defaulted and you owe it money on the swaps. What are your options?
\end{interviewq}

\begin{interviewq}[$\star\star$ \iqroles{trader, risk}]\label{iq:fm:legal-documentation:5}
Why does close-out netting matter for the capital a bank holds against a counterparty?
\end{interviewq}

\begin{interviewq}[$\star\star\star$ \iqroles{risk, researcher}]\label{iq:fm:legal-documentation:6}
How would you estimate the probability that a fund trips at least one \omterm{def:fm:legal-documentation:ate}{NAV trigger} in the next year?
\end{interviewq}
